\documentclass[10pt,a4paper]{amsart}
\usepackage[margin=19mm]{geometry}
\usepackage{amsmath,amssymb,mathtools}
\usepackage[scr=rsfs,cal=euler]{mathalfa}
\usepackage{xfrac}
\usepackage[unicode,hidelinks,bookmarks=false]{hyperref}
\newcommand{\ie}{i.e.\ }
\newcommand{\cf}{cf.\ }
\newcommand{\resp}{respectively\ }
\newcommand{\Subsection}{\subsection}
\providecommand{\nobreakdash}{\nobreak\hbox{-}}
\setlength{\emergencystretch}{2em}
\multlinegap0pt
\usepackage[shortlabels]{enumitem}
\usepackage{amssymb}
\usepackage{xcolor}
\usepackage{tikz}
\usepackage{tcolorbox}
\newtheorem{lemma}{Lemma}
\newtheorem{theorem}{Theorem}
\newcommand{\II}{{\mathcal I}}
\newcommand{\Matrixc}[1]{\ensuremath{\left[\begin{array}{cccccccccccc|ccc} #1 \end{array}\right]}}
\newcommand{\PH}{H}
\title{Automated Classification of Homeostasis Structure in Input-Output Networks}
\author{Xinni Lin, Fernando Antoneli,, Yangyang Wang}
\date{Source: arXiv:2603.08882v1 (CC BY 4.0)}
\begin{document}
\maketitle

\noindent\emph{Source and licence.} Automated Classification of Homeostasis Structure in Input-Output Networks. Version arXiv:2603.08882v1 (CC BY 4.0), distributed by arXiv under the Creative Commons Attribution 4.0 International licence, http://creativecommons.org/licenses/by/4.0/, as recorded in the arXiv licence metadata for this exact version and shown on the abstract page https://arxiv.org/abs/2603.08882v1. Full licence notice: \texttt{licenses/arxiv-2603.08882-CC-BY-4.0.txt}.\par\smallskip

\noindent\textit{Editorial scope.} Counted unit: the article's theorem thm:GenH (source0.tex lines 630--637). The article's complete original proof of it is delivered with it (source0.tex lines 639--702, one \texttt{proof} environment of 64 lines). No mathematics is written, repaired or re-derived: the statement and the proof are the article's own lines, in the article's own notation, and no retained line is substituted or re-flowed.  Retained uncounted as the source-local context the proof needs: lines 238--246; lines 288--298; lines 570--582.  Nothing was omitted from any retained span, so the package carries no omission record.  No retained line contains a citation, so the wrapper carries no bibliography and no bibliography entry.  No retained line contains a figure environment, an image-inclusion command or drawing code, so no image is delivered and the package declares no asset dependency.  Editorial formatting only, in addition: an AMS \texttt{amsart} preamble that restates the article's own declarations and macros used by the retained lines, namely the environments \texttt{lemma}, \texttt{theorem} on separate counters, together with the standard packages those lines need.  The retained environments print in this document as Lemma 1, Theorem 1 in order of appearance, whereas the article prints the same items with the numbering of its own full document.  The complete native source of the article as distributed in the arXiv e-print for this exact version is bundled unchanged as \texttt{sources/196035/source0.tex}.
\begin{equation} \label{eq: io}
\begin{split}
\dot{x}_\iota & = 
f_\iota(x_\iota,x_\rho, x_o,\II) \\
\dot{x}_\rho & =
f_\rho(x_\iota,x_\rho,x_o) \\
\dot{x}_o & = f_o(x_\iota,x_\rho,x_o)
\end{split}
\end{equation}

\begin{lemma} \label{L:det}
Let $(X_0,\II_0)$ be an asymptotically stable equilibrium of \eqref{eq: io}.  The input-output function $x_o(\II)$ satisfies
\begin{equation} \label{xo'_new}
x_o' = \pm\frac{f_{\iota, \II}}{\det(J)} \det(\PH)
\end{equation}
Hence, $\II_0$ is a point of infinitesimal homeostasis if and only if 
\begin{equation} \label{xo'_reduced}
\det(\PH)= 0
\end{equation}
at $(X_o, \II_0)$.
\end{lemma}

\begin{equation} \label{weighted_homeostasis_matrix_definition}
\langle H \rangle = 
\begin{bmatrix}
f_{\iota, x_\iota} &  f_{\iota, x_\rho} & -f_{\iota, \mathcal{I}} \\
f_{\rho, x_\iota}&  f_{\rho, x_\rho} & 0 \\
f_{o, x_\iota} &  f_{o, x_\rho} & 0
\end{bmatrix} = \begin{bmatrix} f_{\iota_{1}, x_{\iota_{1}}} & \cdots & f_{\iota_{1}, x_{\iota_{n}}} & f_{\iota_{1}, x_{\rho}} & - f_{\iota_{1}, \mathcal{I}} \\
    \vdots & \ddots & \vdots & \vdots & \vdots \\
    f_{\iota_{n}, x_{\iota_{1}}} & \cdots & f_{\iota_{n}, x_{\iota_{n}}} & f_{\iota_{n}, x_{\rho}} & - f_{\iota_{n}, \mathcal{I}} \\
    f_{\rho, x_{\iota_{1}}} & \cdots & f_{\rho, x_{\iota_{n}}} & f_{\rho, x_{\rho}} & 0 \\
    f_{o, x_{\iota_{1}}} & \cdots & f_{o, x_{\iota_{n}}} & f_{o, x_{\rho}} & 0
    \end{bmatrix}
\end{equation}

\begin{theorem} \label{thm:GenH}
Let $\mathcal{G}$ be a single-input input-output network with $n$ distinguished input nodes 
$\iota = (\iota_1, \dots, \iota_n)$ receiving the same input from $\mathcal{I}\in \mathbb{R}$, $N$ regulatory nodes $\rho$, and one distinguished output node $o$.
Construct the \emph{augmented network} $\mathcal{G}^\diamond$ by introducing a new input node $I$ and converting the original input nodes $\iota = (\iota_1, \dots, \iota_n)$ into regulatory nodes, by adding directed edges 
$I \to \iota_j$ for $j = 1,\dots,n$. 
Extend the vector field $F=(f_{\iota},f_{\rho},f_o)$ to $F^\diamond=(f_{I},f_{\iota},f_{\rho},f_o)$ in such a way that the equation for the new input node variable  $\dot{x}_I(t)=f_I(x_{I},\mathcal{I})$ enforces $x_I(t)\equiv \II$ asymptotically. 
Then the infinitesimal homeostasis matrix $H^\diamond$ of $\mathcal{G}^\diamond$ is identical to the generalized homeostasis matrix $\langle H \rangle$ of $\mathcal{G}$ given by \eqref{weighted_homeostasis_matrix_definition}, up to a sign difference.
\end{theorem}

\begin{proof} 
% The generalized homeostasis matrix \eqref{weighted_homeostasis_matrix_definition} of $\mathcal{G}$ can be rewritten as
% $$
% \langle H \rangle = \Matrixc{
% f_{\iota, x_{\iota}} & f_{\iota, x_{\rho}} & -f_{\iota, \mathcal{I}}\\
% f_{\rho, x_{\iota}} & f_{\rho, x_{\rho}} & 0\\
% f_{o, x_{\iota}} & f_{o, x_{\rho}} & 0\\
% }
% $$
Let $H^\diamond$ be the homeostasis matrix for the augmented network $\mathcal{G}^\diamond$. 
We prove below that $\det(H^\diamond) = \det(\langle H \rangle)$, as polynomials in $(f_{\ell,x_j})$, up to a sign.
Consequently, both formulations yield identical infinitesimal homeostasis conditions.
We define the dynamics of the new input node as 
\[
\dot{x}_{I}=f_I(x_{I},\mathcal{I}) := \mathcal{I}-x_I,
\]
which ensures $x_I(t) = \mathcal{I}$ asymptotically.
Moreover, because of the skew product from of $F^\diamond$, if the original system $F$ is stable at an equilibrium $(X_0,\II_0)$ then the extended system $F^\diamond$ is stable at the corresponding equilibrium $(\II_0,X_0,\II_0)$. 
Other choices of dynamics are possible. 
This is one working example. 
We then rewrite dynamics for the original set of input nodes as 
\[
\dot{x}_{\iota} = f_{\iota}(x_\iota,x_\rho,x_o, x_I),
\]
which now become regulatory nodes that do not directly depend on the input parameter $\mathcal{I}$. 
The Jacobian matrix $J^\diamond$ of $\mathcal{G}^\diamond$ becomes:
\[
J^\diamond = \begin{bmatrix}
f_{I, x_{I}} & f_{I, x_{\iota}} & f_{I, x_{\rho}} & f_{I, x_o} \\
f_{\iota, x_{I}} & f_{\iota, x_{\iota}} & f_{\iota, x_{\rho}} & f_{\iota, x_o} \\
f_{\rho, x_{I}} & f_{\rho, x_{\iota}} & f_{\rho, x_{\rho}} & f_{\rho, x_o} \\
f_{o, x_{I}} & f_{o, x_{\iota}} & f_{o, x_{\rho}} & f_{o, x_o}
\end{bmatrix}
\quad=\quad
\begin{bmatrix}
f_{I, x_{I}} & 0 & 0 & 0 \\
f_{\iota, x_{I}} & f_{\iota, x_{\iota}} & f_{\iota, x_{\rho}} & f_{\iota, x_o} \\
0 & f_{\rho, x_{\iota}} & f_{\rho, x_{\rho}} & f_{\rho, x_o} \\
0 & f_{o, x_{\iota}} & f_{o, x_{\rho}} & f_{o, x_o}
\end{bmatrix}
\]
By Lemma \ref{L:det}, the input-output function $x_o(\mathcal{I})$ of this augmented network satisfies 
\[
x_o'=\pm \frac{f_{I,\mathcal{I}}}{\det(J^\diamond)}\det(H^\diamond) = \pm \frac{\det(H^\diamond)}{\det(J)}
\]
where $H^\diamond$ is obtained from removing the first row and the last column of $J^\diamond$:
\begin{equation}
\begin{aligned}
H^\diamond &= \begin{bmatrix}
f_{\iota, x_{I}} & f_{\iota, x_{\iota}} & f_{\iota, x_{\rho}}\\
0 & f_{\rho, x_{\iota}} & f_{\rho, x_{\rho}}\\
0 & f_{o, x_{\iota}} & f_{o, x_{\rho}}
\end{bmatrix} 
\quad = \quad
\begin{bmatrix}
f_{\iota, \mathcal{I}} & f_{\iota, x_{\iota}} & f_{\iota, x_{\rho}}\\
0 & f_{\rho, x_{\iota}} & f_{\rho, x_{\rho}}\\
0 & f_{o, x_{\iota}} & f_{o, x_{\rho}}
\end{bmatrix}
\end{aligned}
\end{equation}
The second equality follows that, at the stable equilibrium, $x_I=\mathcal{I}$. 
Through two column exchanges, we obtain that $\det(H^\diamond) = \det(\langle H \rangle)$. 
\end{proof}

\end{document}
